\chapter{Data, ethics and external validity}
\label{app:ethics}

\section{What the data is, and what it reveals}

The corpus consists of anonymised signalling records collected by a Czech mobile
operator over fifteen days in March 2014 and provided to the university for
research. There are no names, no phone numbers and no content of any kind: a row
holds a per-day pseudonymous identifier, a small amount of demographic metadata,
and a list of (cell, timestamp) pairs.

That is nonetheless sensitive material, and it is worth being explicit about why.
A day of cell attachments is a coarse but genuine trace of where a person was and
when. Applied to one individual it can reveal where they sleep, where they spend
working hours, when they leave and return, and whether their routine changed on a
given day --- which is precisely the information the ``home cell'' and ``activity
cell'' classification in this corpus extracts. The fact that a record is
generated automatically, without the person doing anything, is what makes the
trace continuous and therefore informative.

\section{How the work handled it}

Four properties of the method limit exposure, and three of them were chosen for
methodological reasons that happen to align with data minimisation.

\begin{itemize}
  \item \textbf{The analysis never leaves cell or site granularity.} No attempt
        was made to interpolate positions between records, to reconstruct a route,
        or to attach a cell to an address or a point of interest. The finest
        spatial unit anywhere in this report is one of 369 radio cells over 113
        base stations, and mobility is measured on the coarser of the two.
  \item \textbf{Time is reduced to the hour.} The raw timestamps are
        second-resolution; every model input in this report uses the hour of day
        only. That was decided because twenty-four values are learnable and
        86{,}400 are not, but it also removes the fine timing that makes a trace
        individually identifying.
  \item \textbf{Only the trajectory columns are used.} Age, gender and the
        operator's own behaviour class were never fed to any model.
  \item \textbf{Identifiers are per day and are not reconciled.} A person appears
        under a different identifier in each day-file, and no attempt was made to
        link them --- which is also, as Chapter~\ref{ch:scale} explains, the
        limitation that bounds what the work could conclude.
  \item \textbf{Nothing left the research environment.} The data stayed on the
        university's files and on the working copies used for the experiments; no
        raw record appears in this report, and the only trajectories quoted are
        short fragments used as illustrations.
\end{itemize}

\section{External validity}

Three limits should be attached to any use of these results outside their
setting.

\paragraph{One area, one operator, one fortnight.} Everything here comes from a
single Czech region observed over fifteen days in March 2014, served by 369
cells. Vocabulary coverage of the test set is already complete at five thousand
training users, and the saturation result of Chapter~\ref{ch:scale} partly
reflects how small that world is. A dense urban network with orders of magnitude
more cells would not necessarily saturate at the same point.

\paragraph{Single-day results.} Every number in this report is the end of one day
predicted from its beginning. Nothing here supports a claim about predicting one
day from previous days, and the experiment intended to probe day-of-week
structure did not complete.

\paragraph{Day-files contain disjoint user populations.} Because no user appears
in two day-files, a change of day-file is a change of test population, and
therefore a change of baseline. Comparisons across day-files are only valid as
deltas to each run's own baseline; the absolute accuracies are not comparable.
This applies to the two headline results of the last iterations, which were
measured on two different 100-user test sets whose baselines differ by 5.1
points.
